%% 第 20 章：正规算子
%% 本文件由 tools/md2tex.py 自动生成，请勿直接编辑。

\chapter{正规算子}


\begin{jdefn}{正规算子的定义}{r:21:1}
若：

\[
T^*T=TT^*,
\]

则称 $T$ 是正规算子。

自伴算子一定正规，因为：

\[
T^*T=TT=TT^*.
\]

但正规算子不一定自伴。平面旋转算子就是正规但非自伴的例子。
\end{jdefn}

\begin{jprop}{正规算子的范数刻画}{r:21:2}
\[
T\text{ 正规}
\iff
\|Tv\|=\|T^*v\|
\text{ 对所有 }v\in V.
\]
\end{jprop}

\begin{jproof}{证明思路}
将两边范数平方展开：

\[
\|Tv\|^2
=
\langle T^*Tv,v\rangle,
\]

\[
\|T^*v\|^2
=
\langle TT^*v,v\rangle.
\]

因此，若 $T^*T=TT^*$，两者范数相同。

反过来，若所有向量的两个范数都相同，在复内积空间中，这足以推出：

\[
T^*T=TT^*.
\]
\end{jproof}

\begin{jprop}{正规算子的特征向量也是伴随的特征向量}{r:21:3}
若 $T$ 正规，且：

\[
Tv=\lambda v,
\qquad
v\ne0,
\]

则：

\[
T^*v=\overline{\lambda}v.
\]
\end{jprop}

\begin{jproof}{证明思路}
算子：

\[
T-\lambda I
\]

仍然正规。

而：

\[
(T-\lambda I)v=0.
\]

正规算子的范数刻画说明：若一个向量被正规算子映为零，它也会被该算子的伴随映为零。

因此：

\[
(T-\lambda I)^*v=0.
\]

又因为：

\[
(T-\lambda I)^*
=
T^*-\overline{\lambda}I,
\]

所以：

\[
T^*v=\overline{\lambda}v.
\]
\end{jproof}

\begin{jprop}{不同特征值对应的特征向量正交}{r:21:4}
若 $T$ 正规，且：

\[
Tv=\lambda v,
\qquad
Tw=\mu w,
\qquad
\lambda\ne\mu,
\]

则：

\[
\langle v,w\rangle=0.
\]
\end{jprop}

\begin{jproof}{证明思路}
由前一结论：

\[
T^*w=\overline{\mu}w.
\]

于是：

\[
\langle Tv,w\rangle
=
\lambda\langle v,w\rangle.
\]

另一方面：

\[
\langle Tv,w\rangle
=
\langle v,T^*w\rangle
=
\mu\langle v,w\rangle.
\]

因此：

\[
(\lambda-\mu)\langle v,w\rangle=0.
\]

因为：

\[
\lambda\ne\mu,
\]

所以：

\[
\langle v,w\rangle=0.
\]
\end{jproof}
